% Ad Atticum 14.11 (ad-atticum-14-11)
% Source: https://raw.githubusercontent.com/PerseusDL/canonical-latinLit/master/data/phi0474/phi057/phi0474.phi057.perseus-lat1.xml
% Fetched by scripts/fetch_latin.py

% Dateline (Perseus): Scr. in Cumano xi K. Mai. a. 710 (44).

\ciceroLetterOpener{CICERO ATTICO salutem}

\ciceroSection{1} nudius tertius dedi ad te epistulam longiorem; nunc ad ea quae proxime. velim me hercule Asturae Brutus. ἀκολασίαν istorum scribis. an censebas aliter? equidem etiam maiora exspecto. quom equidem contionem lego de tanto viro, de clarissimo civi, ferre non queo. etsi ista iam ad risum. sed memento, sic alitur consuetudo perditarum contionum ut nostri illi non heroes sed di futuri quidem in gloria sempiterna sint, sed non sine invidia, ne sine periculo quidem. verum illis magna consolatio conscientia maximi et clarissimi facti, nobis quae? qui interfecto rege liberi non sumus. sed haec fortuna viderit, quoniam ratio non gubernat.

\ciceroSection{2} de Cicerone quae scribis iucunda mihi sunt; velim sint prospera. quod curae tibi est ut ei suppeditetur ad usum et cultum copiose per mihi gratum est, idque ut facias te etiam atque etiam rogo. de Buthrotiis et tu recte cogitas et ego non dimitto istam curam. suscipiam omnem etiam actionem quam video cotidie faciliorem. de Cluviano, quoniam in re mea me ipsum diligentia vincis, res ad centena perducitur. ruina rem non fecit deteriorem, haud scio an etiam fructuosiorem. hic mecum Balbus, Hirtius, Pansa. modo venit Octavius et quidem in proximam villam Philippi mihi totus deditus; Lentulus Spinther hodie apud me; cras mane vadit.
